\section{A category of analytic operators with explicit bounds}
\label{ct:operators}

The earlier completion and category packets already supply a general
extension principle and elementary category laws. The useful addition here
is a small operator calculus that records which analytic operations descend
to a completion, how their error bounds compose, and which apparently
similar operations do not descend. This section is a specialization and
quantitative synthesis of familiar mathematics, rather than a claim that
bounded-operator extension is a new classical theorem.

\subsection{Objects, arrows, and rapid names}

A \emph{rational seminormed presentation} consists of a decidable set of
codes, primitive recursive codes for zero, addition, negation and rational
scalar multiplication, and a rational-valued seminorm satisfying the vector
space and seminorm laws. The equality of codes may be finer than equality
of represented vectors; zero seminorm identifies the redundant vectors.
All stated laws must have object-theory proofs or occur as explicit
hypotheses. Dyadic step functions with their $L^1$ seminorm are the principal
example. Their refinement relations ensure that equivalent finite
presentations have zero distance.

Write $e_n=2^{-n}$. A rapid name $a$ satisfies
\[
   \|a_q-a_p\|\leq e_p\qquad(q\geq p).
\]
The represented equality of two names is the cross estimate
\[
 a\simeq b\quad\Longleftrightarrow\quad
 \forall p,q\;\|a_p-b_q\|\leq e_p+e_q.
\]
A rapid name satisfies this estimate with itself. For transitivity, insert
$b_j$ into the triangle inequality and obtain the desired bound plus
$2e_j$. A strictly positive rational violation would exceed $2e_j$ for an
explicit sufficiently large $j$, found from its numerator and denominator.
This proves transitivity intuitionistically using rational order. No test
for equality of completed objects has been introduced.

An arrow $T:V\to W$ is a specified primitive recursive linear map, together
with an integer $M\geq0$ and a proof of
\[
                    \|Tv\|\leq M\|v\|.
\]
Equality of arrows means equality of their values in the target
seminormed quotient. The same underlying arrow can have many valid bounds.
Thus the bound is operational data, not part of extensional arrow equality.
Define
\[
 h(M)=\min\{j\leq M+1:M\leq2^j\}.
\]
The set is nonempty, and this is bounded search. In particular $h(0)=h(1)=0$.

\begin{theorem}[Quantitative bounded-operator extension]
\label{ct:extension}
The formula
\[
                   (\widehat T a)_n=T(a_{n+h(M)})
\]
constructs a rapid name. It preserves represented equality, agrees with $T$
on constant names, and defines a linear map on the completion. Identity and
composition are preserved up to represented equality. If $\|Sv\|\leq N\|v\|$,
the bound attached to $S\circ T$ is $NM$.
\end{theorem}
\begin{proof}
For $q\geq p$, linearity and the norm bound give
\[
 \|T(a_{q+h(M)})-T(a_{p+h(M)})\|
 \leq M e_{p+h(M)}\leq e_p.
\]
For equivalent inputs the corresponding cross estimate is at most
$M(e_{p+h(M)}+e_{q+h(M)})\leq e_p+e_q$.
Addition of names is represented by $a_{n+1}+b_{n+1}$; its rapidness follows
by adding the two input estimates. Scalar multiplication uses the same
construction with an integer upper bound for the absolute scalar.
Linearity of the extended map follows by applying the finite linearity
identity to sufficiently accurate approximations and then removing their
rational error bounds. The resulting finite argument uses no induction on
a formula quantifying over names.

The two constructors for a composite evaluate $ST$ at indices
$n+h(NM)$ and $n+h(N)+h(M)$. Rapidness bounds the distance between these
input values by the larger of their two errors, hence also by their sum.
After multiplication by $NM$, each sum term is at most $e_n$.
Using arbitrary output indices $p,q$ therefore gives the equality bound
$e_p+e_q$. Identity has shift zero. This proves the category laws on
represented arrows directly. An extension is unique among bounded linear
extensions: approximate the input by $a_n$, use the common core value, and
bound the two remaining differences by the respective operator bounds
times $e_n$.\end{proof}

This proof explains an efficiency gain. An integration, averaging, or
conditional-expectation operator with norm at most one needs no extra
precision shift. A pipeline $T_r\cdots T_1$ can use a single shift for the
product of the bounds instead of repeatedly rebuilding names. If separately
implemented stages introduce errors $\delta_i$, the output error is at most
\[
 \sum_{i=1}^r\delta_i\prod_{j=i+1}^rM_j.
\]
The proof is finite induction on the stages using
$E_{i+1}\leq M_{i+1}E_i+\delta_{i+1}$. The invariant is a rational
inequality after denominators are cleared. This is an object-iPRA induction
candidate with a quantifier-free invariant, rather than an assertion about
an unbounded family of analytic operators.

\subsection{Bilinear operations and the correct tensor product}

\begin{theorem}[Bilinear extension on certified bounded inputs]
\label{ct:bilinear}
Let a specified rational bilinear map satisfy
$\|B(u,v)\|\leq M\|u\|\|v\|$.
Suppose rapid names $a,b$ have uniform integer bounds
$\|a_n\|\leq R$, $\|b_n\|\leq S$.
For $H=h(M(R+S))$, the sequence
\[
                       B(a_{n+H},b_{n+H})
\]
is rapid. It depends only on the represented inputs, and it is bilinear on
completion classes.
\end{theorem}
\begin{proof}
The finite identity
\[
 B(u,v)-B(u',v')=B(u-u',v)+B(u',v-v')
\]
gives
\[
 \|B(u,v)-B(u',v')\|
 \leq M\bigl(\|u-u'\|\|v\|+\|u'\|\|v-v'\|\bigr).
\]
For two indices $q\geq p$ the right side is at most
$M(R+S)e_{p+H}\leq e_p$. To compare two constructions with different certified bounds, first enlarge
both pairs of bounds to their coordinatewise maxima, and use the resulting
common shift. For a fixed input pair, comparing its old and enlarged shifts
uses its old bound $M(R+S)$; both shifts dominate that bound, so the cross
error is at most $e_p+e_q$. With the common bounds and common shift in place,
the cross-name inequalities prove equality for equivalent input pairs.
Transitivity therefore proves independence of both names and bound choices.
The finite bilinear identities pass through these same estimates.
Every rapid name admits such an integer bound: use
$R=\lceil\|a_0\|\rceil+2$, since
$\|a_n\|\leq\|a_0\|+1$.\end{proof}

For step functions on separate probability coordinates, define
$(f\otimes g)(x,y)=f(x)g(y)$. Finite rectangular summation gives exactly
\[
               \|f\otimes g\|_1=\|f\|_1\|g\|_1.
\]
The theorem therefore supplies a well-defined tensor operation
$L^1(X)\times L^1(Y)\to L^1(X\times Y)$ on the represented completions.
The same finite factorization gives
$I(f\otimes g)=I(f)I(g)$; finite sums of tensors and the Fubini contraction
from the measure section extend this identity to the corresponding named
objects. Associativity, symmetry, and the unit are witnessed by reindexing
finite rectangular arrays, then by the extension estimates.

This claim concerns this concrete tensor operation. It does not silently
construct the projective norm on an arbitrary algebraic tensor product:
that norm is an infimum over decompositions, and an effective approximation
to that infimum would require further data.

Multiplication \emph{on the same coordinate} is different. If
$f_m=2^{m/2}\mathbf1_{[0,2^{-m})}$ for even $m$, then
$\|f_m\|_1=2^{-m/2}$ but $\|f_m^2\|_1=1$. Thus multiplication cannot be a
continuous bilinear map $L^1\times L^1\to L^1$. It becomes bounded on
$L^\infty\times L^1$ when an explicit uniform bound is supplied, or on
$L^2\times L^2\to L^1$ with a Cauchy--Schwarz certificate. The typing of an
analytic expression must distinguish these cases.

\subsection{Probability kernels as contraction arrows}

Let $p_i,q_j$ be rational probability vectors and let $K_{ij}\geq0$ satisfy
\[
       \sum_jK_{ij}=1,\qquad q_j=\sum_i p_iK_{ij}.
\]
For a target vector $f$ define $(K^*f)_i=\sum_jK_{ij}f_j$.
Then
\[
 \sum_i p_i|(K^*f)_i|
 \leq\sum_{i,j}p_iK_{ij}|f_j|
 =\sum_jq_j|f_j|.
\]
This is a contraction in the weighted $L^1$ seminorms. If $L$ is a second
kernel, ordinary finite multiplication gives
$(KL)^*=K^*L^*$ and the required marginal identity. Positivity and the
constant-one vector are preserved. These identities give a category of
rational probability presentations and law-compatible kernels, acting
contravariantly by contractions on observables.

To pass a family of finite kernels to completed probability presentations,
provide compatibility with refinement, or an explicit summable operator
error bound for that compatibility. Existence of an arbitrary measurable
conditional probability is not inferred from this construction. The
specific dyadic averaging kernels used for Fubini and conditional
expectation have exact refinement compatibility, so the construction is
unconditional for those kernels.

\section{Integral-quantifier semantics and its substitution discipline}
\label{ct:logic}

Integration can act as a real-valued logical quantifier; this is an
established model-theoretic direction, represented for example by Bagheri
and Pourmahdian's continuous integration logic
\cite{ct-integration-logic}. Here the goal is a much smaller numerical
semantics whose every demanded value has an explicit finite approximation.
The construction makes no appeal to compactness, ultraproducts, or the
existence of models of arbitrary consistent theories.

\subsection{Syntax with measured contexts}

For $d$ free variables use the dyadic probability cube
$\Omega_d=[0,1)^d$, with $\Omega_0$ a one-point probability space. Rational
step functions are stored on finite product partitions. A formula in
context $d$ is generated by:
\begin{enumerate}
\item an atomic predicate occurrence with a supplied rapid $L^1(\Omega_d)$
name whose finite values lie in $[0,1]$;
\item a specified rational connective $c:[0,1]^r\to[0,1]$, with rational
coefficients $\lambda_i\geq0$ satisfying
$|c(u)-c(v)|\leq\sum_i\lambda_i|u_i-v_i|$;
\item a quantifier $Q\psi$ in context $d$ when $\psi$ has context $d+1$,
interpreted by integration in the last variable;
\item a pullback along a supplied map of measured contexts with the
bounded-compression certificate described below.
\end{enumerate}
One can replace rational sensitivity coefficients by integer upper bounds
when a simple shift is desired. Product, complement, minimum, maximum,
truncated subtraction, and rational convex combinations all admit such
bounds. All constituent approximations are refined to a common finite
partition before a connective is evaluated.

The qualification ``atomic predicate occurrence'' is deliberate: an
$L^1$ name for a two-variable predicate cannot automatically be reused as
a name for its diagonal occurrence $P(x,x)$. That restriction is proved,
and a useful sufficient repair is given, below.

\begin{theorem}[Compositional $L^1$ sensitivity]
\label{ct:formula}
Give each atomic occurrence $\alpha$ an error $\epsilon_\alpha$ between
two interpretations in its actual measured context. Define the error
expression recursively by
\[
 E_\alpha=\epsilon_\alpha,\quad
 E_{c(\psi_1,\ldots,\psi_r)}=\sum_i\lambda_iE_{\psi_i},\quad
 E_{Q\psi}=E_\psi,\quad E_{\tau^*\psi}=D_\tau E_\psi.
\]
Then the $L^1$ distance between the two formula interpretations is at most
$E_\varphi$. In particular, a sentence has absolute value error at most
$E_\varphi$. If all atomic names are rapid, all operations are exact on
finite step functions, and $W_\varphi$ is obtained by first replacing each rational sensitivity
and compression coefficient by its nonnegative integer ceiling and then
setting every atomic error to one, evaluating all atoms at
$n+h(W_\varphi)$ gives a rapid name for the value of $\varphi$.
\end{theorem}
\begin{proof}
At a connective, integrate its pointwise sensitivity inequality and use
finite linearity of the integral. At $Q$, the finite inequality
\[
 \left|\sum_j w_j(f_{ij}-g_{ij})\right|
 \leq\sum_jw_j|f_{ij}-g_{ij}|
\]
is summed over the remaining coordinates. This is precisely the Fubini
contraction, and it survives passage to the rapid completions. The pullback
case is the contraction estimate with coefficient $D_\tau$ below. These
arguments prove the error recurrence by induction on a finite formula.
For two precision indices $q\geq p$, the atomic error is at most $e_p$;
therefore the whole error is at most $W_\varphi e_p$. The specified shift
makes it at most $e_n$. A sentence is an element of the completion of
$\mathbb Q$ since its measured context has one point.\end{proof}

This is a form of approximate elementary preservation for the chosen
fragment: if two presented structures have small atomic discrepancies in
all the contexts actually used, the displayed error expression bounds the
change in every fixed formula. It is not a completeness theorem for a
logical proof system. For a proof-producing interpretation, the finite
sensitivity lemmas are first derived in the arithmetic calculus, and every
rule of the selected integral logic must then be checked against them.

There is a practical gain over uniform-net evaluation: integration does not
increase the sensitivity coefficient. Nested expectation quantifiers can
be evaluated by successive exact averaging. Fubini allows their order to
be chosen to minimize intermediate table size. The formula proof justifies
this change of order. The representation may still have exponentially
many cells; no polynomial complexity bound is asserted.

\subsection{Which substitutions are legal in \texorpdfstring{$L^1$}{L1}?}

A step-compatible map $\tau:X\to Y$ has \emph{compression bound} $D$ if,
for each target partition used and every target cell $A_j$,
\[
                \mu_X(\tau^{-1}(A_j))\leq D\mu_Y(A_j).
\]
The preimages must have specified finite descriptions or certified named
approximations, with their own error budget. In the exact finite case,
collecting the source cell weights proves
\[
 \|g\circ\tau\|_{L^1(X)}
 =\sum_j|g_j|\mu_X(\tau^{-1}(A_j))
 \leq D\|g\|_{L^1(Y)}.
\]
Thus pullback extends by Theorem~\ref{ct:extension}. Compression bounds
multiply under composition. Coordinate permutations and projections with
independent unused coordinates have $D=1$. They justify variable renaming
and the appropriate weakening operations.

\begin{proposition}[Diagonal substitution and supremum are not $L^1$ rules]
\label{ct:diagonal}
There is no continuous operator on arbitrary $L^1([0,1)^2)$ classes that
agrees with $f\mapsto(x\mapsto f(x,x))$ on all dyadic step functions.
There is also no continuous $L^1$ operation agreeing with supremum on all
bounded dyadic step functions.
\end{proposition}
\begin{proof}
Let $A_{m,i}=[i2^{-m},(i+1)2^{-m})$ and set
\[
 u_m(x,y)=\sum_{i<2^m}\mathbf1_{A_{m,i}}(x)
                              \mathbf1_{A_{m,i}}(y).
\]
The $2^m$ occupied squares each have area $2^{-2m}$, so
$\|u_m\|_1=2^{-m}$. Yet $u_m(x,x)=1$ for every $x$.
The input sequence converges to zero in $L^1$, while the proposed output
sequence stays distance one from zero. In fact the diagonal map has no
finite compression bound: the union of the occupied squares has area
$2^{-m}$ and its diagonal preimage has measure one.
For supremum use $v_m=\mathbf1_{[0,2^{-m})}$, which has $L^1$ norm
$2^{-m}$ and supremum one.\end{proof}

Consequently one must not treat arbitrary almost-everywhere predicate
classes as a full first-order structure with unrestricted substitutions.
A safe implementation either supplies a separate atomic name in each
required context, limits substitutions to certified compression maps, or
upgrades the predicates to uniformly controlled pointwise representatives.
Exact zero tests and sharp truth thresholds likewise require separation;
compositional approximation alone cannot decide them.

\subsection{Recovering pointwise predicates from quantitative regularity}

There is a useful repair that is stronger than merely forbidding the
problematic substitutions. Work on the closed cube $[0,1]^d$; the boundary
has no effect on the integral. Suppose $f,g$ have specified pointwise
representations and are both $L$-Lipschitz for the maximum metric.
All integrals in the following proof can be constructed by finite rational
Riemann sums with the supplied Lipschitz modulus.

\begin{theorem}[$L^1$-to-uniform upgrade]
\label{ct:upgrade}
For $d\geq1$ and $0<\delta\leq1$,
\[
 \sup_x|f(x)-g(x)|
 \leq \delta^{-d}\|f-g\|_1+2L\delta.
 \tag{\ref{ct:upgrade}a}
\]
Let $(f_n)$ be a family with this common Lipschitz bound and a supplied
$L^1$ Cauchy modulus $N$, so that
$\|f_n-f_m\|_1\leq2^{-r}$ for $n,m\geq N(r)$.
Put
\[
 a=k+h(8L+1),\qquad r=k+da+1.
\]
Then $N(r)$ is a uniform Cauchy index for error $2^{-k}$.
Here an integer upper bound for $L$ is used in $h$.
\end{theorem}
\begin{proof}
Fix $x$. In each coordinate choose an interval of length $\delta$ contained
in $[0,1]$ and containing $x_i$; for example its left endpoint can be
$\min(x_i,1-\delta)$. Their product $C_x$ has volume $\delta^d$, and
$\|x-y\|_\infty\leq\delta$ for every $y\in C_x$. Hence
\[
 |f(x)-g(x)|\leq|f(y)-g(y)|+2L\delta.
\]
Integrating this inequality on $C_x$, bounding its integral by the
whole-cube integral, and dividing by the positive rational $\delta^d$
gives the displayed estimate. When $x$ is rational this is a direct limit
of finite rational inequalities. For a represented $x$, approximate it by
rational points and use the common Lipschitz bound; its value estimate
passes to the limit with an explicitly chosen input precision.

Take $\delta=2^{-a}$. Since $2^{h(8L+1)}\geq8L+1$, the second term is at
most $2^{-k}/4$. With $r=k+da+1$, the first term is at most
$2^{da}2^{-r}=2^{-k}/2$. Their sum is strictly below $2^{-k}$.
Thus the reindexed family is uniformly Cauchy. A diagonal of its specified
value approximations constructs the pointwise limit, and passing the
finite Lipschitz inequalities to that limit proves the same bound $L$.
\end{proof}

This theorem supplies a primitive recursive modulus upgrade in every fixed
dimension. It requires the common regularity bound and an $L^1$ convergence
modulus; it does not extract either from their mere classical existence.
Once the upgrade has been made, diagonal substitution is continuous for
the supremum norm, so the desired pointwise formula can be evaluated. The
result is also a concrete interface between the present integral semantics
and the earlier uniform finite-net semantics.

\section{Stable topological observations without exact distance tests}
\label{ct:topology}

Finite complexes here are certificates for statements about metric spaces
and their approximations. The target is multiscale topological information,
not the replacement of topology by an unrelated finite problem.
Rips interleavings from metric correspondences are standard in persistence
theory; see Chazal, de Silva, and Oudot \cite{ct-rips}. The contribution here
is the explicit arithmetic certificate, the treatment of uncertain
distances, and the separation of the finite homotopy proof from geometric
reconstruction claims.

\subsection{Rational complexes and the finite homotopy witness}

For a finite rational metric table $(A,d_A)$ and rational scale $r\geq0$,
let $\operatorname{VR}_r(A)$ contain the nonempty finite subsets whose
pairwise distances are at most $r$. A $q$-chain over $\mathbb F_2$ is a
finite parity sum of $(q+1)$-element simplices. Its boundary is the parity
sum of the $q$-element faces. Each twice-deleted face appears exactly twice,
so $\partial^2=0$. In degree zero use the unaugmented boundary zero.

For an ordered vertex list, write $[v_0,\ldots,v_q]$ for the sorted simplex
when the vertices are distinct and for zero otherwise. This normalization
commutes with the boundary over $\mathbb F_2$: in a repeated-vertex list,
faces deleting the two repeated occurrences cancel, while the other faces
remain degenerate. A simplicial vertex map thus induces a chain map even
when it identifies vertices.

Two vertex maps $f,g:K\to L$ are contiguous if
$f(\sigma)\cup g(\sigma)$ is a simplex of $L$ for every simplex $\sigma$
of $K$. Define
\[
 P_q[v_0,\ldots,v_q]
   =\sum_{i=0}^{q}[f(v_0),\ldots,f(v_i),g(v_i),\ldots,g(v_q)].
\]
All its summands are legal by contiguity.

\begin{lemma}[Explicit homotopy for contiguous maps]
\label{ct:prism}
The finite operator above satisfies
$\partial P+P\partial=f_\#+g_\#$.
In particular, for every cycle $z$, the chain $Pz$ witnesses that
$f_\#z$ and $g_\#z$ differ by a boundary.
\end{lemma}
\begin{proof}
Expand the boundary of the $i$th summand. Deleting $f(v_j)$ for $j<i$
cancels the summand of $P\partial$ on the source face omitting $v_j$.
Deleting $g(v_j)$ for $j>i$ cancels its corresponding source-face summand.
The two remaining faces are
\[
 A_i=[f(v_0),\ldots,f(v_{i-1}),g(v_i),\ldots,g(v_q)],
 \qquad A_{i+1}.
\]
As $i$ runs from zero to $q$, $A_1,\ldots,A_q$ occur twice and cancel.
Only $A_0=g_\#[v_0,\ldots,v_q]$ and
$A_{q+1}=f_\#[v_0,\ldots,v_q]$ survive. The normalization observation
above covers degeneracies. For a cycle $\partial z=0$, the identity reduces
to $\partial(Pz)=f_\#z+g_\#z$.\end{proof}

This lemma already appears in the repository's topological work. It is
included to make the new interleaving certificate self-contained; it is
not counted as a new mathematical result of this package.

\subsection{A correspondence produces a certified interleaving}

Let $R\subseteq A\times B$ project onto both finite vertex sets and satisfy
\[
 |d_A(a,a')-d_B(b,b')|\leq\eta
       \quad((a,b),(a',b')\in R).
 \label{ct:distortion}
\]
Choose $f(a)$ and $g(b)$ by least-index search in the nonempty finite rows
and columns of $R$. This is bounded search, not an infinite choice rule.

\begin{theorem}[Finite distortion certificate]
\label{ct:rips}
The maps $f,g$ induce chain maps at shifted scales
\[
 F_r:C_*(\operatorname{VR}_r(A))\longrightarrow
           C_*(\operatorname{VR}_{r+\eta}(B)),\qquad
 G_r:C_*(\operatorname{VR}_r(B))\longrightarrow
           C_*(\operatorname{VR}_{r+\eta}(A)).
\]
Their composites are explicitly chain homotopic to the inclusion maps at
shift $2\eta$. Different selections subordinate to $R$ yield chain-homotopic
maps already at shift $\eta$. The induced homology interleaving is therefore
independent of the finite tie-breaking convention.
\end{theorem}
\begin{proof}
If $a,a'$ lie in a scale-$r$ simplex, the distortion bound gives
$d_B(f(a),f(a'))\leq r+\eta$, proving simpliciality.
For $h=gf$, the relations $(a,f(a)),(h(a),f(a))\in R$ give
$d_A(a,h(a))\leq\eta$. More sharply, for vertices $a,a'$ in that simplex,
\[
 d_A(a,h(a'))\leq d_B(f(a),f(a'))+\eta\leq r+2\eta,
\]
and the same bound holds for $d_A(h(a),h(a'))$ by applying the distortion
estimate to $(h(a),f(a))$ and $(h(a'),f(a'))$.
Thus the original simplex and its $h$-image together have diameter at most
$r+2\eta$. Lemma~\ref{ct:prism} supplies the required chain homotopy.
The calculation for $fg$ is identical with $A,B$ exchanged.

If $f,f'$ are two selections, every mixed pair in
$f(\sigma)\cup f'(\sigma)$ is a pair of $R$-images of source vertices.
Its distance is at most $r+\eta$. Hence the maps are contiguous at that
scale. Finally the vertex maps do not change when the scale increases,
so all maps commute with the scale-inclusion maps. This is exactly the
interleaving naturality equation.\end{proof}

A chain with $s$ supported $q$-simplices produces at most $s(q+1)$ prism
summands before parity cancellation. A relation on $n$ and $m$ vertices
requires at most $n^2m^2$ rational distortion comparisons in a naive
checker. These are explicit primitive recursive bounds. The homology
statement can remain the equation
$\partial(Pz)=GFz+z$ for each finite cycle; construction of an internal
quotient object is unnecessary for this endpoint.

\subsection{Uncertain real distances need scale slack, not sign decisions}

Suppose $x_1,\ldots,x_n$ are represented points and a rational symmetric
table $\widetilde d_{ij}$ satisfies
$|\widetilde d_{ij}-d(x_i,x_j)|\leq\eta$.
The vertices are the labelled indices $1,\ldots,n$, and the true distances
form a pseudometric on those indices. Coincident represented points are
not identified by an equality test. Write $\underline x$ for this labelled
presentation. The approximate table need not itself satisfy the triangle
inequality. Define the finite clique complex $K_r$ by the decidable tests
$\widetilde d_{ij}\leq r$. Directly from the certificate,
\[
 \operatorname{VR}_{r-\eta}(\underline x)
       \subseteq K_r
       \subseteq\operatorname{VR}_{r+\eta}(\underline x),
 \qquad r\geq\eta.
 \label{ct:sandwich}
\]
If a second approximation has error $\eta'$, then
$|\widetilde d_{ij}-\widetilde d'_{ij}|\leq\eta+\eta'$, so the identity on
vertices induces
\[
             K_r\subseteq K'_{r+\eta+\eta'},\qquad
             K'_r\subseteq K_{r+\eta+\eta'}.
\]
The composites are literally the inclusion maps. This gives a stable
multiscale answer even when an exact edge distance equals the threshold.
It does not decide the homology of the exact complex at that threshold.
Requiring a gap for every distance comparison would unnecessarily exclude
this useful stable observation.

\subsection{From a compact metric presentation to finite topology}

Fix a rational $\epsilon>0$. Let $X$ be a rational metric core with a finite $\epsilon$-net $A\subseteq X$
and a supplied primitive recursive selector $s:X\to A$ satisfying
$d(x,s(x))\leq\epsilon$. Inclusion $j:A\to X$ and the selector induce
\[
 s_r:\operatorname{VR}_r(X)\to\operatorname{VR}_{r+2\epsilon}(A),
 \qquad j_r:\operatorname{VR}_r(A)\to\operatorname{VR}_r(X).
\]
For a source simplex, all distances between its vertices and the selected
vertices are at most $r+2\epsilon$. Thus $js$ is contiguous to inclusion
at shift $2\epsilon$; the same argument proves that $sj$ is contiguous to
inclusion on $A$. One obtains an asymmetric interleaving with forward
shift $2\epsilon$ and backward shift zero.

In particular the scale-inclusion homology map
\[
 H_q(\operatorname{VR}_r(X))\longrightarrow
 H_q(\operatorname{VR}_{r+2\epsilon}(X))
\]
factors through the homology of a finite complex on $|A|$ vertices. Its
rank is therefore at most
$\binom{|A|}{q+1}$, or the simpler primitive recursive bound
$|A|^{q+1}$. Internally the factorization is again the explicit finite-chain
identity provided by the prism, with the net selector applied only to the
vertices in the input chain. It does not require quantification over a
completed infinite homology vector space.

For completion points, take a sufficiently accurate core approximation,
then apply a finer core net selector. For example a core net of radius
$\epsilon/2$ and a core approximation of error at most $\epsilon/2$ give
an $\epsilon$-net selection on the represented points. A finite chain needs
only finitely many such value calls. A fixed primitive recursive family
supplies them directly; arbitrary program-code names require an explicit
clocked batch or an existing proof of batch totality.

What this proves is persistent topological control of a supplied metric
presentation. It does not identify a small-scale Rips complex with the
singular homology of an arbitrary space. Such identification requires
additional geometric hypotheses and a reconstruction argument. No
manifold, reach, local-contractibility, or good-cover bound is manufactured
by the finite checker.

\subsection{Exact arithmetic endpoints}

The main topology endpoint is the finite implication
\[
 \mathrm{Metric}(A)\wedge\mathrm{Metric}(B)\wedge
 \mathrm{CorrDist}(R,\eta)\wedge\mathrm{Cycle}_q(z,r)
 \;\Longrightarrow\;
 \partial\,\mathrm{Prism}(g f,\mathrm{id},z)=GFz+z.
\]
All finite-table conditions can be Booleanized; the polynomial/rational
comparisons use positive denominator clearing. The loop invariant records
the processed coefficient parity at each output simplex. Formula
sensitivity has the analogous endpoint consisting of a finite-step
$L^1$ error inequality. The bounded-operator extension uses two selected
name indices and their supplied inequalities. These are precise targets
for object-iPRA derivations. The human proofs above establish the
mathematics and the finite induction structure; no new kernel-checked
\texttt{Prf PRAAx} proof is claimed by this section alone.
